% AY2016 制御工学 第2問〔II〕（転記）
% 出典: 03_制御工学/original/03_制御工学/2016/2016se.pdf
% 要約: G1,G2,G3 の直列に2つの負帰還（G2 まわり、および G1 後段ノードから入力側）。
%       伝達関数・ゲイン・ボード線図・ステップ応答・過渡特性の比較考察。

\section*{第2問〔II〕}

下図のシステムについて次の問い (1)〜(4) に答えよ. 入力を $U$, 出力を $Y$ とする.

\begin{center}
\begin{tikzpicture}[>=Latex]
  \node (in)   at (0,0)   {$U$};
  \node (sum1) [sum] at (1.2,0) {};
  \node (G1)   [block] at (2.7,0) {$G_1$};
  \node (sum2) [sum] at (4.3,0) {};
  \coordinate (nA) at (5.1,0);
  \node (G2)   [block] at (6.3,0) {$G_2$};
  \coordinate (nB) at (7.5,0);
  \node (G3)   [block] at (8.7,0) {$G_3$};
  \node (out)  at (10.1,0) {$Y$};

  \draw[->] (in)   -- (sum1);
  \draw[->] (sum1) -- (G1);
  \draw[->] (G1)   -- (sum2);
  \draw[->] (sum2) -- (G2);
  \draw[->] (G2)   -- (G3);
  \draw[->] (G3)   -- (out);
  \fill (nA) circle (1.4pt);
  \fill (nB) circle (1.4pt);

  % 内側負帰還: nB → sum2（上から）
  \draw[->] (nB) -- (7.5,1.3) -- (4.3,1.3) -- (sum2.north);
  % 外側負帰還: nA → sum1（下から）
  \draw[->] (nA) -- (5.1,-1.3) -- (1.2,-1.3) -- (sum1.south);

  % 符号
  \node at ($(sum1.north west)+(0.02,0.10)$) {\scriptsize $+$};
  \node at ($(sum1.south east)+(-0.02,-0.10)$) {\scriptsize $-$};
  \node at ($(sum2.north east)+(-0.02,0.10)$) {\scriptsize $-$};
  \node at ($(sum2.south west)+(0.02,-0.10)$) {\scriptsize $+$};
\end{tikzpicture}
\end{center}

\begin{enumerate}[label=(\arabic*)]
  \item 伝達関数を求めよ.

  \item $G_1=\dfrac{3}{s}$, $G_2=\dfrac{1}{s+4}$, $G_3=4$ のときのゲイン $R$ を求めよ.
        また, ゲイン線図を描け.

  \item (2) のステップ応答を求め, 時間応答概形を描け.

  \item $G_3=2s+4$ のときのステップ応答を求めよ. (2) と比較し, ゲイン線図ならびに
        時間応答の両観点から過渡特性について考察せよ.
\end{enumerate}
